\documentclass[aspectratio=169]{beamer}

\usepackage{listings}
\usepackage{listings-rust}
\usepackage{xcolor, tikz}
\usepackage{linehighlight}
\usepackage{fontspec}
\usepackage{wrapfig}
\usepackage{animate}
\usetikzlibrary{shapes,arrows, positioning, decorations.markings}

\usetheme[sectionpage=simple]{metropolis}

\title{An Overview of the Embedded Rust Ecosystem}
\subtitle{Or, Why are there so many crates, and what do they do?}
\author{Frans Skarman aka. TheZoq2}


\definecolor{background}{rgb}{0.98,0.98,0.98}
\definecolor{oxidizepurple}{RGB}{49 0 86}

\setbeamercovered{invisible}
\setbeamercovered{%
  again covered={\opaqueness<1->{50}}}
\setbeamercolor{title separator}{bg=background, fg=oxidizepurple}
\setbeamercolor{frametitle}{bg=background, fg=oxidizepurple}


\setmainfont{Ubuntu}
% \usepackage{fontspec}
% \setmonofont[Contextuals={Alternate}]{Hasklug Nerd Font}

% set colors for code background and highlighting
\definecolor{codehighlight}{rgb}{0.95,0.9,0.95}
\definecolor{codebackground}{rgb}{0.98,0.98,0.98}
\definecolor{mygreen}{rgb}{0,0.4,0}
\definecolor{mygray}{rgb}{0.5,0.5,0.5}
\definecolor{mymauve}{rgb}{0.58,0,0.82}
\lstset{ 
    % choose the background color; you must add \usepackage{color} or \usepackage{xcolor}; should come as last argument
    % backgroundcolor=\color{white},
    % the size of the fonts that are used for the code
    basicstyle=\footnotesize\ttfamily,
    commentstyle=\color{mygreen},
    % sets if automatic breaks should only happen at whitespace
    breakatwhitespace=false,
    % sets automatic line breaking
    breaklines=true,
    % if you want to delete keywords from the given language
    deletekeywords={...},
    % if you want to add LaTeX within your code
    escapeinside={\%*}{*)},
    % lets you use non-ASCII characters; for 8-bits encodings only, does not work with UTF-8
    extendedchars=true,
    % keeps spaces in text, useful for keeping indentation of code (possibly needs columns=flexible)
    keepspaces=true,
    % keyword style
    keywordstyle=\color{blue},
    % where to put the line-numbers; possible values are (none, left, right)
    numbers=none,
    % how far the line-numbers are from the code
    numbersep=5pt,
    % the style that is used for the line-numbers
    %numberstyle=\ttfamily\footnotesize\color{mygray},
    % if not set, the frame-color may be changed on line-breaks within not-black text (e.g. comments (green here))
    rulecolor=\color{black},
    % show spaces everywhere adding particular underscores; it overrides 'showstringspaces'
    showspaces=false,
    % underline spaces within strings only
    showstringspaces=false,
    % show tabs within strings adding particular underscores
    showtabs=false,
    % the step between two line-numbers. If it's 1, each line will be numbered
    stepnumber=1,
    % string literal style
    stringstyle=\color{mymauve},
    % sets default tabsize to 2 spaces
    tabsize=2,
}


\begin{document}
\maketitle

% \begin{frame}{TODO}
%     \begin{itemize}
%         \item Ferris attribution
%     \end{itemize}
% \end{frame}

\begin{frame}
    \frametitle{Who am I?}
    
    Frans Skarman (@TheZoq2)

    \begin{itemize}
        \item Rust evangelist
        \item Maintainer of \texttt{stm32f1xx\_hal}
        \item My desk is covered in rust
    \end{itemize}

    % \animategraphics[autoplay, loop, width=0.4\textwidth]{30}{fig/wave/}{0}{30}
    \includegraphics[width=0.5\textwidth]{fig/wave.png}
\end{frame}

\begin{frame}
    \frametitle{Who are you?}

    \begin{itemize}[<+>]
        \item Embedded rust
        \item Arduino
        \item Embedded C
        \item Rust user
    \end{itemize}
\end{frame}

\begin{frame}{Goals}
    \uncover<1>{There are lots of embedded crates}

    \uncover<1>{How do they all fit together?}

    \uncover<2>{Explained through showing the need for them}

    \uncover<3>{The initial code is not something you'd write 😉}
\end{frame}

\begin{frame}
    \frametitle{The lowest level}

    \begin{itemize}
        \item Micro controllers control I/O through ``registers''
        \item Typically memory mapped
        \item Read the datasheet for instructions
    \end{itemize}
\end{frame}

\begin{frame}{An example: turning on an LED}
    \uncover<1>{Using an \texttt{stm32f1} processor}

    \uncover<2>{
        Start by reading the 1000 page datasheet:
    }

    \uncover<3>{
        Process:
        \begin{itemize}
            \item{Power up the GPIO peripheral in the RCC peripheral}
            \item{Configure pin as output}
            \item{In the correct mode}
            \item{Set output to High}
        \end{itemize}
    }
\end{frame}

\begin{frame}{An example: turning on an LED}
    \begin{columns}
        \begin{column}{0.4\textwidth}
            Configure pin 13 as output

            \begin{itemize}
                \item{Control register at offset \texttt{0x04}}
                \item{Offset from start of GPIOx}
                \item{Write \texttt{0b10} in bit 20, 21}
                \item{Write \texttt{0b01} in bit 23, 22}
            \end{itemize}

            Same deal for the output value.
        \end{column}
        \begin{column}{0.6\textwidth}
            \begin{figure}
                \includegraphics[width=\textwidth]{fig/gpio_crl.png}
            \end{figure}
        \end{column}
    \end{columns}
\end{frame}

\begin{frame}{The rust code}
    \vspace{-0.5cm}
    \begin{linehighlight}{
          \only<1>{
                \highline{2,...,2}
          }
          \only<2>{
                \highline{4,...,6}
          }
          \only<3>{
                \highline{8,...,9}
          }
          \only<4>{
                \highline{11,...,17}
          }
    }
        \lstinputlisting[language=rust]{code/raw.rs}
    \end{linehighlight}
\end{frame}

\begin{frame}{The rust code}
    \begin{itemize}
        \item Very unsafe
        \item Very unergonomic
    \end{itemize}
    \hspace*{5cm}\includegraphics[width=0.48\textwidth]{fig/unsafe.png}
\end{frame}

\begin{frame}{SVD files}
    \begin{itemize}
        \item{Published by microcontroller manufacturers}
        \item{Describes the function of each register}
        \item{\texttt{svd2rust} generates rust crates}
    \end{itemize}
\end{frame}

\begin{frame}{Peripheral Access Crates (PACs)}
    \begin{itemize}
        \item{Result of patches + \texttt{svd2rust}}
        \item{\textit{Mostly} safe interface}
        \item{Adds a lot of zero cost abstraction}
        \item{Prevents re-use of peripherals}
    \end{itemize}
\end{frame}

\begin{frame}{Peripheral Access Crates (PACs)}
    \begin{linehighlight}{
          \only<1>{
                \highline{2,...,2}
          }
          \only<2>{
                \highline{4,...,9}
          }
          \only<3>{
                \highline{10,...,14}
          }
    }
        \lstinputlisting[language=rust]{code/pac.rs}
    \end{linehighlight}
\end{frame}

\begin{frame}{Peripheral Access Crates (PACs)}
    \begin{itemize}
        \item{Much higher level interface}
        \item{Zero cost abstraction}
        \begin{itemize}\item{(In release mode)}\end{itemize}
        \item{Most unsafety is now fixed}
    \end{itemize}
    \hspace*{5cm}\includegraphics[width=0.40\textwidth]{fig/rustacean-flat-happy.pdf}
\end{frame}

\begin{frame}{Peripheral Access Crates (PACs)}
    Some remaining issues
    \begin{itemize}
        \item{No check for peripheral dependencies}
        \item{Or correct intialisation}
    \end{itemize}

    $\Rightarrow$ Still requires through reading of the datasheet
    \hspace*{5cm}\includegraphics[width=0.40\textwidth]{fig/ferrisSob.png}
\end{frame}

\begin{frame}{Hardware Abstraction Layers (HALs)}
    A high level interface around the PAC
\end{frame}

\begin{frame}[fragile]{Hardware Abstraction Layers (HALs)}

    \begin{linehighlight}{
          \only<1>{ % Only on slides 1 and 3 (beamer stuff)
                \highline{2,...,3} % highlight code lines 1,3 and 5.
          }
          \only<2> {
                \highline{5,...,6} % highlight lines 2 to 8.
          }
          \only<3>{ % Only on slide 2 (beamer)
                \highline{8,...,11} % highlight lines 2 to 8.
          }
    }
        \lstinputlisting[language=rust]{code/hal_example.rs}
          % put your \lstinputlisting{} or \begin{lstlisting}...\end{lstlisting} here
    \end{linehighlight}
\end{frame}

\begin{frame}{Hardware Abstraction Layers (HALs)}
    \begin{itemize}
        \item{High level interface}
        \item{Type state ensures correct initialisation}
    \end{itemize}
    \hspace*{5cm}\includegraphics[width=0.40\textwidth]{fig/rustacean-flat-happy.pdf}
\end{frame}


\begin{frame}{External Hardware}
    Using our HAL to talk to other hardware

    HY-SRF05 ultrasonic distance sensor

    \begin{itemize}
        \item Send a pulse on one pin
        \item Wait for pulse on another
        \item Length of return pulse is proportional to distance
    \end{itemize}

    \hspace*{5cm}\includegraphics[width=0.40\textwidth]{fig/ultrasonic.png}
\end{frame}

\begin{frame}{External Hardware (Implementation)}
    \vspace{-0.5cm}
    \begin{linehighlight}{
          \only<1>{ % Only on slides 1 and 3 (beamer stuff)
                \highline{3,...,4} % highlight code lines 1,3 and 5.
          }
          \only<2>{ % Only on slide 2 (beamer)
                \highline{6,...,9} % highlight lines 2 to 8.
          }
          \only<3>{ % Only on slide 2 (beamer)
                \highline{11,...,16} % highlight lines 2 to 8.
          }
    }
        \lstinputlisting[language=rust]{code/ultrasonic.rs}
    \end{linehighlight}
\end{frame}

\begin{frame}{External Hardware (Usage)}
    \begin{linehighlight}{
          \only<1>{ % Only on slides 1 and 3 (beamer stuff)
                \highline{2,...,3} % highlight code lines 1,3 and 5.
          }
          \only<2>{ % Only on slide 2 (beamer)
                \highline{5,...,9} % highlight lines 2 to 8.
          }
          \only<3>{ % Only on slide 2 (beamer)
                \highline{11,...,12} % highlight lines 0 to 8.
          }
    }
        \lstinputlisting[language=rust]{code/ultrasonic_usage.rs}
    \end{linehighlight}
\end{frame}

\begin{frame}{External Hardware}
    We just made our first ``driver''

    Types ensure that things are configured correctly

    Ownership ensures pins and timers aren't shared
    \hspace*{5cm}\includegraphics[width=0.40\textwidth]{fig/rustacean-flat-happy.pdf}
\end{frame}


\begin{frame}{Another HAL}
    HAL for another microcontroller

    Probably does similar things

    With a slightly different API

    \vspace*{1cm}\hspace*{2cm}\includegraphics[width=0.40\textwidth]{fig/another_hal.pdf}
    % \hspace*{5cm}\includegraphics[width=0.30\textwidth]{fig/rustacean-flat-happy.pdf}
\end{frame}

\begin{frame}{Another HAL}
    \begin{itemize}
        \item Our driver used our types
        \item It is not compatible with the other HAL
    \end{itemize}
    \includegraphics[width=0.4\textwidth]{fig/static_sweat.png}
    % \includegraphics[width=0.4\textwidth]{fig/rustacean-flat-sad.pdf}
\end{frame}


\begin{frame}{Embedded HAL}
    \begin{itemize}
        \item Provides traits for common peripherals
        \begin{itemize}
            \item GPIO
            \item Timers
            \item SPI
            \item UART
            \item etc...
        \end{itemize}
        \item HALs implement them
        \item Drivers use them
        \item Driver re-use across the ecosystem!
        % TODO: Ferrises
    \end{itemize}
\end{frame}

\begin{frame}{HAL Usage Example}
    \lstinputlisting[language=rust]{code/ultrasonic_eh.rs}
\end{frame}

\begin{frame}
    \section{Misc. Crates}
\end{frame}

\begin{frame}{Real-Time Interrupt-driven Concurrency (RTIC)}
    \begin{itemize}
        \item A concurrency framework for building real time systems
        \begin{itemize}
            \item Tasks
            \item Task communication
            \item Real time analysis
            \item etc.
        \end{itemize}
        \item Previously known RTFM (Real Time For the Masses)
        \item{See the previous talk by Emil Fresk}
        \hspace*{7cm}\includegraphics[width=0.2\textwidth]{fig/rtic.png}
    \end{itemize}
\end{frame}

\begin{frame}{Board Support Crates (BSP)}
    Provides utilities for breakout boards

    \begin{itemize}
        \item Silkscreen pinouts
        \item Drivers for peripherals
        \item etc.
    \end{itemize}
\end{frame}

\begin{frame}{\texttt{cortex-x} crates}
    Cortex-<m,a,r> is the core of many arm microcontrollers

    Some relevant crates:
    \begin{itemize}
        \item \texttt{cortex-m}: PAC for core peripherals (timers, debugger etc.)
        \item \texttt{cortex-m-rt}: minimal runtime (entry points, exceptions, \dots)
            \begin{itemize}\item included with \texttt{rt} feature in PACs\end{itemize}
        \item \texttt{cortex-m-semihosting}: Debugging convenience
    \end{itemize}
\end{frame}

\begin{frame}{Conclusion}
    \uncover<1> {
        \begin{itemize}
            \item \texttt{embedded-hal} is the core of ecosystem
            \item HALs implement embedded-hal traits
            \item usually by using PACs
            \item drivers use \texttt{e-h} traits to talk to external devices
        \end{itemize}
    }
    \uncover<2> {
        Other crates exist
        \begin{itemize}
            \item{RTIC provides concurrency}
            \item{svd2rust generates PACs}
        \end{itemize}
    }
\end{frame}

\begin{frame}{Where should you start?}
    \uncover<1> {
        HALs are nice!

        Starting from scratch?
        \begin{itemize}
            \item{Pick a microcontroller with a good HAL}
        \end{itemize}
    }

    \uncover<2>{
        Specific hardware in mind?
        \begin{itemize}
            \item{Look for a HAL}
            \item{If none exist}
            \begin{itemize}
                \item Check for a PAC
                \item Or make one using svd2rust if svd files are available
                \item Build a HAL on top
            \end{itemize}
        \end{itemize}
    }

    \uncover<3> {
        Similar steps for drivers
    }
\end{frame}

\begin{frame}{Resources}
    \begin{itemize}
        \item A collection of most rust-embedded related things
            \begin{itemize}\item \url{https://github.com/rust-embedded/awesome-embedded-rust}\end{itemize}
        \item{embedded-wg chatroom}
        \item{Community discord}
    \end{itemize}

\end{frame}

\begin{frame}
    \section{Thank you for listening!}
\end{frame}

\end{document}
